\documentclass[11pt]{article}
\usepackage[margin=1in]{geometry}
\usepackage[T1]{fontenc}
\usepackage{lmodern}
\usepackage{amsmath,amssymb,amsthm}
\usepackage{graphicx}
\usepackage{microtype}
\usepackage[colorlinks=true,linkcolor=blue,citecolor=blue,urlcolor=blue]{hyperref}

\newtheorem{theorem}{Theorem}
\newtheorem{lemma}[theorem]{Lemma}
\newtheorem{corollary}[theorem]{Corollary}
\newcommand{\R}{\mathbb R}
\newcommand{\ext}{\operatorname{ext}}
\newcommand{\cconv}{\overline{\operatorname{conv}}}

\title{A weak-star-convergent-basis obstruction\\
for a polyhedral Lindenstrauss problem}
\author{Candidate partial result for arXiv:2303.10023}
\date{11 August 2026}

\begin{document}
\maketitle

\begin{center}
\fbox{\parbox{0.91\linewidth}{\textbf{Status.} Candidate partial result,
likely valid, requiring human review. Problem 6.8 remains open in general.
The theorem below rules out every $\ell_1$-predual for which the standard
dual basis is weak-star convergent.}}
\end{center}

\section{Source problem}

De Bernardi proves that every polyhedral Lindenstrauss space admits
arbitrarily close equivalent polyhedral norms whose unit balls are generated
by their extreme points. He then asks \cite[Problem~6.8]{DeBernardi}:
\begin{quote}
Does there exist a polyhedral Lindenstrauss space $E$ satisfying
\[
                 B_E=\cconv\bigl(\ext(B_E)\bigr)?
\]
\end{quote}
As throughout the source paper, ``space'' is understood to be
infinite-dimensional unless otherwise stated.

\begin{figure}[ht]
\centering
\includegraphics[width=\linewidth]{figures/open_problem_crop.png}
\caption{Problem 6.8 and its immediate context in
\cite[p.~20]{DeBernardi}, rendered and cropped from the local source PDF.}
\end{figure}

For $f=(f_n)\in\ell_1$, put
\begin{equation}\label{eq:Wf}
 W_f=\left\{x=(x_n)\in c:
       \lim_{n\to\infty}x_n=\sum_{n=1}^{\infty}f_nx_n\right\},
\end{equation}
with the supremum norm inherited from $c$. Hyperplanes of this form are a
standard source of explicit $\ell_1$-preduals; the polyhedral example in
\cite{CMP} is obtained with $f_n=2^{-n}$.

\section{Proof intuition}

Suppose the standard basis $(e_n^*)$ of $X^*=\ell_1$ converges weak-star to
$f$. Evaluating $x\in X$ on these basis vectors embeds $X$ isometrically into
$c$. The weak-star limit imposes precisely the single equation defining
$W_f$, and an annihilator calculation shows that the range is all of $W_f$.

The extreme-point calculation is then coordinatewise. If two coordinates of
$x\in B_{W_f}$ are strictly inside $[-1,1]$, a small two-coordinate
perturbation can be chosen to preserve the defining equation. Thus an extreme
point has at most one unsaturated coordinate. Its limit must consequently be
$s\in\{-1,1\}$. The equality
\[
 1=|s|=\left|\sum f_nx_n\right|
       \le \sum |f_n||x_n|\le\|f\|_1\le1
\]
forces $\|f\|_1=1$ and $x_n=s\,\operatorname{sgn}(f_n)$ on the support of
$f$; in fact every coordinate is saturated. If the support of $f$ has two
points, all convex combinations of extreme points obey a fixed coordinate
equation that an elementary vector of $B_{W_f}$ violates. If the support has
only one point, $W_f$ is just an isometric copy of the nonpolyhedral space
$c$.

\section{The hyperplane representation}

\begin{lemma}\label{lem:representation}
Let $X$ be a real Banach space and fix an isometric identification
$X^*=\ell_1$. If the standard basis $(e_n^*)$ converges in
$\sigma(\ell_1,X)$ to $f\in\ell_1$, then $\|f\|_1\le1$ and $X$ is linearly
isometric to $W_f$.
\end{lemma}

\begin{proof}
Weak-star lower semicontinuity of the norm gives $\|f\|_1\le1$. Define
\[
 J:X\longrightarrow c,\qquad Jx=(e_n^*(x))_{n=1}^{\infty}.
\]
The sequence is convergent because $e_n^*\to f$ weak-star. Moreover,
$\ext(B_{X^*})=\{\pm e_n^*:n\in\mathbb N\}$ is a boundary of $X$, so
\[
                 \|Jx\|_\infty=\sup_n|e_n^*(x)|=\|x\|.
\]
Thus $J$ is an isometry with closed range $Y=JX$. For $x\in X$,
\[
 \lim_n(Jx)_n=f(x)=\sum_{n=1}^{\infty}f_ne_n^*(x)
                  =\sum_{n=1}^{\infty}f_n(Jx)_n,
\]
so $Y\subseteq W_f$.

Use the standard representation
\[
 c^*=\R\oplus_1\ell_1,
 \qquad (a_0,a)(z)=a_0\lim_nz_n+\sum_na_nz_n.
\]
Such a functional annihilates $Y$ exactly when
\[
            a_0f+a=0\quad\hbox{in }X^*=\ell_1.
\]
Consequently
\[
 Y^\perp=\operatorname{span}\{(1,-f)\}.
\]
But $W_f=\ker(1,-f)$ has the same annihilator. Two closed subspaces with the
same annihilator coincide, hence $Y=W_f$.
\end{proof}

\section{The extreme-point calculation}

Write $\operatorname{supp}f=\{n:f_n\ne0\}$.

\begin{lemma}\label{lem:extreme}
Let $f\in\ell_1$ with $\|f\|_1\le1$. If $\|f\|_1<1$, then
$\ext(B_{W_f})=\varnothing$. If $\|f\|_1=1$, the extreme points are exactly
the sequences $x\in c$ for which there is $s\in\{-1,1\}$ such that
\begin{align}
 x_n&=s\,\operatorname{sgn}(f_n)
       &&(n\in\operatorname{supp}f),\label{eq:on-support}\\
 x_n&\in\{-1,1\} &&(n\in\mathbb N),\label{eq:saturated}\\
 x_n&=s &&\text{eventually}.\label{eq:eventually}
\end{align}
In particular, this set is empty unless $f_n\ge0$ for all sufficiently large
$n\in\operatorname{supp}f$.
\end{lemma}

\begin{proof}
Fix $x\in B_{W_f}$ and let
\[
                     U=\{n:|x_n|<1\}.
\]
If $U$ contains two distinct indices $i,j$, then $x$ is not extreme. Indeed,
if one of $f_i,f_j$ is zero, a sufficiently small nonzero multiple of the
corresponding coordinate vector gives $h\in c_0$ with
$\sum_nf_nh_n=0$ and $x\pm h\in B_c$. If both coefficients are nonzero,
use instead a sufficiently small multiple of
\[
                         f_je_i-f_ie_j.
\]
Again $h\in c_0$, $\sum_nf_nh_n=0$, and $x\pm h\in B_c$. In either case
$h\in W_f$, so $x$ is a nontrivial midpoint in $B_{W_f}$. The same
one-coordinate perturbation shows that if $U=\{j\}$ for an extreme point,
then necessarily $f_j\ne0$.

It follows that an extreme $x$ has at most one unsaturated coordinate. Since
$x$ converges, all but at most one of its coordinates lie in $\{-1,1\}$, and
therefore $\lim_nx_n=s$ for some $s\in\{-1,1\}$ and $x_n=s$ eventually.
The defining equation for $W_f$ now gives
\[
 1=|s|=\left|\sum_nf_nx_n\right|
       \le\sum_n|f_n||x_n|\le\|f\|_1\le1.
\]
Every inequality is an equality. Hence $\|f\|_1=1$, every coordinate with
$f_n\ne0$ is saturated, and all nonzero summands $f_nx_n$ have sign $s$.
Thus \eqref{eq:on-support} holds. The only possible unsaturated coordinate
would now have $f_j=0$, contradicting the necessary condition above; hence
\eqref{eq:saturated} also holds. This proves the necessity of
\eqref{eq:on-support}--\eqref{eq:eventually}, and also proves that there are
no extreme points when $\|f\|_1<1$.

Conversely, a sequence satisfying these three conditions lies in $W_f$:
the equality case gives
$\sum_nf_nx_n=s\sum_n|f_n|=s=\lim_nx_n$. It is an extreme point of the
larger cube $B_c$, because every coordinate is $1$ or $-1$. It is therefore
extreme in the subset $B_{W_f}$ as well.
\end{proof}

\section{The obstruction}

\begin{theorem}\label{thm:main}
Let $X$ be an infinite-dimensional real $\ell_1$-predual for which, under
some isometric identification $X^*=\ell_1$, the standard basis is weak-star
convergent. Then $X$ cannot satisfy both
\begin{enumerate}
\item $X$ is polyhedral, and
\item $B_X=\cconv(\ext(B_X))$.
\end{enumerate}
Consequently no such $X$ answers Problem 6.8 of \cite{DeBernardi}.
\end{theorem}

\begin{proof}
By Lemma~\ref{lem:representation}, $X$ is isometric to $W_f$ for some
$f\in B_{\ell_1}$. If $\|f\|_1<1$, Lemma~\ref{lem:extreme} gives
$\ext(B_{W_f})=\varnothing$, so condition (2) fails.

Suppose $\|f\|_1=1$ and $\operatorname{supp}f$ contains distinct indices
$i,j$. Every extreme point $x$ satisfies
\begin{equation}\label{eq:hyperplane-obstruction}
 \operatorname{sgn}(f_i)x_i=\operatorname{sgn}(f_j)x_j
\end{equation}
by Lemma~\ref{lem:extreme}. Hence the closed convex hull of all extreme
points is contained in the closed hyperplane defined by
\eqref{eq:hyperplane-obstruction}. On the other hand, a sufficiently small
nonzero multiple $y$ of
\[
                         f_je_i-f_ie_j
\]
belongs to $B_{W_f}$, since it has limit zero and $\sum_nf_ny_n=0$.
Its two signed coordinates in \eqref{eq:hyperplane-obstruction} have
opposite signs, so it does not belong to that hyperplane. Thus condition (2)
again fails.

It remains to consider $\operatorname{supp}f=\{k\}$. Then
$f=\sigma e_k$ for $\sigma\in\{-1,1\}$ and
\[
                   W_f=\{x\in c:\lim_nx_n=\sigma x_k\}.
\]
Deleting the $k$th coordinate identifies $W_f$ isometrically with $c$; the
inverse reinserts $\sigma\lim_nx_n$ in that coordinate. The space $c$ is not
polyhedral. For a direct verification, let
\[
 u_n=\cos(1/n),\qquad v_n=\sin(1/n).
\]
Both $u,v$ belong to $c$. On $F=\operatorname{span}\{u,v\}$ the coordinate
functionals correspond to
$p_n=(\cos(1/n),\sin(1/n))\in\R^2$, together with their limit
$p_\infty=(1,0)$. Each $p_n$ is exposed in
\[
       \cconv\{\pm p_n:n\in\mathbb N\}
\]
by the Euclidean functional $p_n$, because its scalar product with every
other $\pm p_m$ and with $\pm p_\infty$ is strictly less than $1$. This set
is the polar of $B_F$, so $B_F$ cannot be a polygon. Hence $c$, and therefore
$W_f$, is not polyhedral. Condition (1) fails in the last case.
\end{proof}

\begin{corollary}
No hyperplane $W_f\subset c$ with $\|f\|_1\le1$ is an
infinite-dimensional polyhedral Lindenstrauss space whose given unit ball is
the closed convex hull of its extreme points. In particular, the explicit
polyhedral hyperplane with $f_n=2^{-n}$ from \cite{CMP} has only the two
extreme points $\pm(1,1,\ldots)$ and cannot answer Problem 6.8.
\end{corollary}

\begin{proof}
The first assertion is the case of Theorem~\ref{thm:main} represented by
Lemma~\ref{lem:representation}; alternatively it follows directly from
Lemmas~\ref{lem:extreme} and the proof of the theorem. For $f_n=2^{-n}$,
the support is all of $\mathbb N$ and every sign is positive, so
Lemma~\ref{lem:extreme} leaves exactly the two constant sign sequences.
\end{proof}

\section{Verification, scope, and search status}

The argument is noncomputational. Its load-bearing points were checked
separately in the accompanying \texttt{verification.md}: the boundary norm
identity in Lemma~\ref{lem:representation}, the annihilator calculation, all
zero-coefficient cases in the perturbation argument, equality in the
$\ell_1$ triangle inequality, and the one-coordinate-support reduction to
$c$.

This is a partial result only. It does not cover an $\ell_1$-predual for
which the dual atoms cannot be enumerated as one weak-star-convergent
sequence, nonseparable $L_1$-preduals, or more general simplex spaces. It
does not answer the source's renorming Problem 6.3 or cardinality Problem
6.10.

\paragraph{Bounded literature check.}
On 11 August 2026, the run indexes and the local full-source corpus were
searched by arXiv id and the terms ``polyhedral Lindenstrauss'', ``hyperplane
of $c$'', ``weak-star convergent basis'', and ``closed convex hull of its
extreme points''. External exact-title, exact-question, DOI, and close-keyword
searches found the source paper and the earlier hyperplane papers
\cite{CMP,CMPV}, but no answer to Problem 6.8 and no statement of
Theorem~\ref{thm:main}. OpenAlex listed no citing work for the 2024 journal
article at the time of the check. This is a bounded novelty audit, not a
claim of definitive priority.

\paragraph{Human-review recommendation.}
Review Lemma~\ref{lem:representation}'s use of the extreme dual boundary and
the equality $Y^\perp=\operatorname{span}(1,-f)$, then check that the
perturbations in Lemma~\ref{lem:extreme} cover every possible unsaturated
coordinate pattern. If those steps pass, the obstruction theorem is ready
to be treated as a rigorous partial result.

\begin{thebibliography}{9}
\bibitem{DeBernardi}
C.~A. De Bernardi,
\emph{Unit balls of polyhedral Banach spaces with many extreme points},
Studia Math. \textbf{275} (2024), 175--196;
arXiv:2303.10023, doi:10.4064/sm230710-31-12.

\bibitem{CMP}
E. Casini, E. Miglierina, and \L. Piasecki,
\emph{A remark on spaces of affine continuous functions on a simplex},
arXiv:1503.09088 (2015).

\bibitem{CMPV}
E. Casini, E. Miglierina, \L. Piasecki, and L. Vesel\'y,
\emph{Rethinking polyhedrality for Lindenstrauss spaces},
Israel J. Math. \textbf{216} (2016), 355--369;
arXiv:1506.08559, doi:10.1007/s11856-016-1412-8.
\end{thebibliography}

\end{document}
